\documentclass{article}
\usepackage{amsmath}
\usepackage{amsfonts}
\usepackage{booktabs}
\usepackage[a4paper, margin=2.5cm]{geometry}
\usepackage{float}
\usepackage{graphicx}
\usepackage{tabularx}
\usepackage[utf8]{inputenc}
\usepackage{longtable}
\usepackage{hyperref}
\usepackage{natbib}
\usepackage{setspace}
\usepackage{array}
\usepackage{amsmath, amsthm}
\setlength{\parskip}{0.45em}
\setlength{\parindent}{0pt}

\usepackage{titlesec}
\titlespacing*{\section}{0pt}{1.0em}{0.5em}

\begin{document}

\title{Referee Report: \\ Identification in Auction Models with Interdependent Costs (Somaini, 2020)}
\author{Jordi Torres}
\date{\today}
\maketitle

\section*{1. Introduction}
This paper studies nonparametric identification of first-price procurement auctions with risk-neutral, asymmetric bidders, dependent private information, and interdependent costs. The motivation is well known: under interdependence, winning gives information about competitor's signals (the winner's curse), and bidders adapt their bids accordingly. The strength of the winner's curse determines how bidders respond to changes in competitive conditions (e.g entry, reserves, subsidies), so being able to quantify the degree of interdependence directly from observed bids is fundamental for substantive counterfactual analysis. Yet, most empirical work has had to pick one of two extreme paradigms ex ante---either pure private values or pure common values---because \citet{LaffontVuong1996} showed that, without further variation, the joint distribution of bids cannot distinguish between these two extreme models.

The author provides a positive identification result for the general interdependent-cost model. He shows that, when bidder-specific cost shifters $x_i$ (e.g., distance from each bidder's plant to the project site) are available and satisfy a reasonable exclusion restriction, the joint distribution of signals and each bidder's full-information cost function can be nonparametrically identified. Identification uses first-order conditions and the rich variation in the set of competitors signals that ``tie for the lowest bid'' induced by continuous shifters. The result is applied to Michigan Department of Transportation highway procurements, where the author estimates the degree of interdependence using a Gaussian parametric structure. Then he uses the identified primitives of the model to evaluate participation- and reserve-price policies.



\section*{2. Model and Main Theoretical Results}

There are $n$ risk-neutral bidders. Bidder $i$ has a private one-dimensional signal $S_i$, observes a public vector of cost shifters $x = (x_1,\dots,x_n)$, and faces ex post completion cost $C_i$. The primitives of the Bayesian game are the joint distribution $F_{C,S\mid x}$ and the resulting full-information cost (FIC):
\[
c_i(s_{-i}, s_i, x_i) \;\equiv\; E\!\left[C_i \,\middle|\, S_{-i}=s_{-i},\, S_i=s_i,\, x\right].
\]
The author imposes the following key assumptions: (A1) signals are independent of cost shifters; (A2) signals are one-dimensional; (A3) bidder $i$'s cost shifter affects only $c_i$ (exclusion); (A4) affiliated signals; (A5) $c_i$ is strictly increasing in $s_i$ and either private or strictly increasing in $s_{-i}$; (A6) shifters generate continuous, sufficiently rich variation; (A7) observables are generated by repeated play of a monotone pure-strategy equilibrium $\beta$ that is invariant across $x$.

To set notation simply, let $\beta_i(s_i, x)$ be bidder $i$'s equilibrium bid function and $\beta_i^{-1}(\cdot, x)$ its inverse, so $s_{-i}(b)\equiv\{\beta_j^{-1}(b,x)\}_{j\neq i}$ is the vector of competitor's signals that make them bid $b$. Let $H_{B_i\mid x}$ denote the conditional CDF of $i$'s equilibrium bid (which is observed). Finally, let $\pi_j^{(i)}$ denote the probability that $i$ ties specifically with bidder $j$ for the lowest bid, conditional on tying with at least one competitor. The main object of interest is bidder $i$'s first-order condition, which after rearranging \citep[eq.~5]{Somaini2020} becomes: 
\[
\underbrace{b \;-\; \frac{\Pr(S_{-i}\ge s_{-i}(b)\mid s_i)}{|\partial_m \Pr(S_{-i}\ge s_{-i}(m)\mid s_i)|_{m=b}|}}_{\xi_i(b,x):\ \text{markup-adjusted bid, identified from bids}}
\;=\; \sum_{j\neq i} \pi_j^{(i)} \, E\!\left[C_i \,\middle|\, s_i, s_j, S_{-ij}\ge s_{-ij}(b), x_i\right].
\]

Intuitively, the FOC states that bidder $i$'s marginal revenue (raising the bid increases the price but lowers the probability of winning) equals the expected FIC conditional on the event of tying with at least one competitor for the lowest bid

The left-hand side $\xi_i$ is identified from the observed bid distribution \citep[as in][]{GuerrePerrigneVuong2000,CampoPerrigneVuong2003}. The right-hand side is a weighted average of expected costs over the events ``tie with $j$, beat all others.'' In the space of signals, each such event ``$S_j=s_j$ and $S_k\ge s_k$ for $k\neq i,j$'' is a face of a rectangle, so their union (the tying event) is an L-shape set when $n=3$ and a higher-dimensional analagous set for $n>3$.

First, identification of the distribution of private values given bids is straight-forward. By A7, monotonicity of bid functions implies $H_{B_i\mid x}=\beta_i^{-1}(\cdot,x)$, so the copula of bids identifies the copula of signals\footnote{ Marginal distributions of signals are normalized to be uniform without loss of generality.The only information that signals carry is ordinal.}

Second, identification of the FIC is a bit more involved and is where his contribution mostly lies. With $n=2$, a single tie set is enough: $\xi_1(b_1,x)=c_1(H_{B_2\mid x}(b_1), H_{B_1\mid x}(b_1), x_1)$. Evaluating at different $(x_2, b_1)$ traces $c_1$ over signal pairs while $x_1$ is held fixed (A3 ensures $x_2$ has no direct effect on $c_1$)\footnote{Note that if we did not have variation from these instruments, we would have a system of equations with more unknowns than equations and thus non-identified. This will become clear with $n>2$}
With $n>2$ the event ``tie for the lowest bid'' is an L-shaped union of $(n-1)$ events, so a single bid yields one equation in $n-1$ unknown expectations. This can't bring identification. The author shows that precisely by varying $x_{-i}$ (which is a vector of size $(n-1)$) moves the corner of the L-shaped set inside the signal space, and that, under certain conditions, one can build an ever-finer partition that recovers $E[C_i\mid s_i, S_{-i}\ge s_{-i}]$ and hence $c_i$ on the relevant set. This leads to the main result of the paper:

\textbf{Theorem 4} states identification of $c_i(s_{-i},s_i,x_i)$ for all $s$ above a threshold, under A1--A3 and A7. The main two conditions that are imposed are: first that observables exhibit sufficient variation to trace down L set; second, they require nonvanishing competition\footnote{Propositions 5--6 establish that monotone equilibria of the first-price auction typically generate observables satisfying these conditions.}

\textbf{Limitations.} The identification result is elegant and meaningfully expands the toolkit researchers can use to study auctions, but several conceptual concerns deserve attention:

First, A7 imposes both equilibrium selection and invariance of the selection across cost shifters. \citet{RenyZamir2004} only guarantee existence of monotone pure-strategy equilibria, not uniqueness. The identification argument requires that, as $x$ varies, the data are generated by a single deterministic rule from $x$ to a particular equilibrium $\beta$. If selection is non-trivial (e.g. bidders coordinate on different equilibria depending on $x$), the copula-based identification of signals is biased and the FIC step inherits this misspecification. 

Second, A2 (one-dimensional signal) is restrictive in the application. Bidders in HMA procurement plausibly receive heterogeneous information about own costs (e.g plant capacity) and about common conditions (input prices, future renegotiation prospects, as the author himself stresses). A2 collapses these into a scalar, which mechanically forces the model to attribute residual cross-bidder dependence to interdependence rather than to a common second signal. Since the empirical contribution rests on distinguishing private from common components, the sensitivity of $\alpha_{i2}$ to violations of A2 is concerning.

Finally, A3 (excluded shifter) is a bit delicate. Distance from $i$'s plant to the site is plausibly excluded from $c_j$ only if plant locations are exogenous. But firms strategically choose plant locations to win in geographic submarkets, and a competitor being close may proxy for unobserved local conditions (local input markets, road density---which is included only crudely as a covariate). The author conditions on auction location, which absorbs some of this concern, but it does not fully address it: a placebo test exploiting variation in distance plausibly orthogonal to local determinants of cost would help.

\section*{3. Empirical Application}
The author estimates the model on 1{,}925 MDOT HMA paving auctions (2001--2010) for 19 frequently observed firms (plus a fringe). The estimation imposes a Gaussian information structure with a two-factor correlation matrix, parametrized so that interdependence is summarized by a single coefficient $\alpha_{i2}$ per bidder. Estimation uses an IV quantile regression with competitors' distance as instrument, following the multistep procedure that follows the identification argument.

\textbf{Key findings.} (i) The data reject pure private values for most bidders but lie strictly between IPC and pure common values: a one-standard-deviation shock to own signal moves FIC by $\approx 10$pp, while the same shock to a competitor's signal moves it by $\approx 4$pp on average and up to $7$pp (Table 7). (ii) Counterfactual restrictions on participation reduce procurement costs, but the decomposition (Table 8) attributes the lion's share of cost savings to the \emph{sampling} effect (selecting the most efficient bidders among invitees, $\approx 80\%$), with the competitive effect ($\approx 9$--$11$pp) offset substantially by the winner's curse ($\approx 35$--$40$pp). (iii) Setting an auction-specific reserve price marginally improves MDOT surplus relative to a binding uniform reserve.

\textbf{Limitations.}
First, the counterfactual on participation policies treats the set of invited bidders as exogenous and equilibrium objects (signal distribution, plant locations, information acquisition) as fixed. In practice, restricting participation likely changes incentives to acquire information \citep{LevinSmith1994} and to locate plants, and the same primitives that determine the winner's curse may shift. This is a classic structural model critique, but here it seems first-order.

Second, the standard errors for the interdependence coefficient are sizable for several bidders (Tables 5--6: ``Common costs'' coefficients often have standard errors of the same order as the point estimate or larger, e.g., bidders 6, 8, 12). The classification of bidders into ``rejects private values'' is largely driven by some frequent participants. The headline counterfactuals could be re-stated with explicit uncertainty bands derived from the bootstrap distribution of $\alpha_{i2}$, to give the reader a sense of how sharp the rejection of IPC is for policy purposes and how general that is.

\section*{4. Conclusion and Further Research}
This is a clear contribution. It closes a long-standing identification puzzle in empirical auctions by showing that, when researchers have access to continuous bidder-specific cost shifters, the interdependent-cost model is positively identified from bids alone---without committing ex ante to either the two extreme paradigms. The argument is constructive, intuitive, and connects naturally with the GPV/CPV tradition: the bidder's FOC delivers expected cost on a tying set, and excluded shifters move the tying set around enough to back out the cost function. The empirical application is honest about the parametric requirements needed for finite-sample estimation, and the counterfactual decomposition (competitive vs.\ affiliation vs.\ winner's curse vs.\ sampling) is clean.

A natural direction for further work is to relax A2 by allowing low-dimensional multidimensional signals (e.g., one private and one common-information signal), which would make the framework directly applicable to settings where bidders have richer information structures. A second direction is to formalize equilibrium-selection robustness: developing a partial-identification version of Theorem 4 under multiple equilibria (or under random equilibrium selection across $x$) would broaden the credibility of the empirical conclusions. Finally, embedding the framework in a model with endogenous participation and information acquisition would make the participation counterfactuals less partial-equilibrium and would also make CF more robust to Lucas type Critiques.

\bibliographystyle{apalike}
\begin{thebibliography}{9}
\bibitem[Campo et~al.(2003)]{CampoPerrigneVuong2003} Campo, S., Perrigne, I., and Vuong, Q. (2003). Asymmetry in first-price auctions with affiliated private values. \emph{Journal of Applied Econometrics}, 18(2), 179--207.
\bibitem[Guerre et~al.(2000)]{GuerrePerrigneVuong2000} Guerre, E., Perrigne, I., and Vuong, Q. (2000). Optimal nonparametric estimation of first-price auctions. \emph{Econometrica}, 68(3), 525--574.
\bibitem[Laffont and Vuong(1996)]{LaffontVuong1996} Laffont, J.-J., and Vuong, Q. (1996). Structural analysis of auction data. \emph{American Economic Review}, 86(2), 414--420.
\bibitem[Levin and Smith(1994)]{LevinSmith1994} Levin, D., and Smith, J. L. (1994). Equilibrium in auctions with entry. \emph{American Economic Review}, 84(3), 585--599.
\bibitem[Reny and Zamir(2004)]{RenyZamir2004} Reny, P. J., and Zamir, S. (2004). On the existence of pure strategy monotone equilibria in asymmetric first-price auctions. \emph{Econometrica}, 72(4), 1105--1125.
\bibitem[Somaini(2020)]{Somaini2020} Somaini, P. (2020). Identification in auction models with interdependent costs. \emph{Journal of Political Economy}, 128(10), 3820--3871.
\end{thebibliography}

\end{document}
